\subsection{CHANGE OF BASE FIELD}

Let \(K_{1}\) be a commutative extension of K. For a Lie algebra g over K to be semi-simple, it is necessary and sufficient that \(g_{(K_{1})}\) be semi-simple; for the Killing form \(\beta_{1}\) of \(g_{(K_{1})}\) is derived from the Killing form \(\beta\) of g by extending the base field from K to \(K_{1}\) ; hence \(\beta_{1}\) is non-degenerate if and only if \(\beta\) is non-degenerate (Algebra, Chapter IX, § 1, no. 4, Corollary to Proposition 3).

If \(\mathfrak{g}_{(\mathbf{K}_{1})}\) is simple, g is semi-simple by the above and cannot be a product of two non-zero ideals, hence g is simple. On the other hand if g is simple \(\mathfrak{g}_{(\mathbf{K}_{1})}\) (which is semi-simple) may be not simple (Exercises 17 and 26 (b)).

Let \(\mathfrak{g}\) be a Lie algebra and \(\mathfrak{r}\) its radical. Then \(\mathfrak{r}_{(\mathbf{K}_1)}\) is the radical of \(\mathfrak{g}_{(\mathbf{K}_1)}\) (§ 5, no. 6). Therefore, if \(\mathfrak{s}\) denotes the nilpotent radical of \(\mathfrak{g}\) , the nilpotent radical of \(\mathfrak{g}_{(\mathbf{K}_1)}\) is \([\mathfrak{g}_{(\mathbf{K}_1)},\mathfrak{r}_{(\mathbf{K}_1)}] = [\mathfrak{g},\mathfrak{r}]_{(\mathbf{K}_1)} = \mathfrak{s}_{(\mathbf{K}_1)}\) . It follows that \(\mathfrak{g}\) is reductive if and only if \(\mathfrak{g}_{(\mathbf{K}_1)}\) is reductive.

Let \(\mathfrak{g}\) be a Lie algebra and \(\mathfrak{h}\) a subalgebra. Recall that a representation of \(\mathfrak{h}\) is semi-simple if and only if the representation of \(\mathfrak{h}_{(\mathrm{K}_1)}\) derived by extending the base field to \(\mathrm{K}_1\) is semi-simple. Hence \(\mathfrak{h}\) is reductive in \(\mathfrak{g}\) if and only if \(\mathfrak{h}_{(\mathrm{K}_1)}\) is reductive in \(\mathfrak{g}_{(\mathrm{K}_1)}\) .

Now let \(K_{0}\) be a subfield of K such that \([K:K_{0}]\) is finite. Let g be a Lie algebra and \(g_{0}\) the (finite-dimensional) Lie algebra derived from g by restricting the base field from K to \(K_{0}\) . Every commutative ideal of g is a commutative ideal of \(g_{0}\) ; conversely, if \(a_{0}\) is a commutative ideal of \(g_{0}\) the smallest vector subspace over K of g containing \(a_{0}\) is a commutative ideal of g; hence g is semi-simple if and only if \(g_{0}\) is semi-simple. If \(g_{0}\) is simple, clearly g is simple. Conversely, suppose that g is simple. We show that \(g_{0}\) is simple. Let \(a_{0}\) be a simple component of \(g_{0}\) . For all \(\lambda \in K^{*}\) , \(\lambda a_{0}\) is an ideal of \(g_{0}\) and

\[
\left[ a _ {0}, \lambda a _ {0} \right] = \lambda \left[ a _ {0}, a _ {0} \right] = \lambda a _ {0} \neq \{0 \},
\]

hence \(\lambda a_0 \supset a_0\) and therefore \(\lambda a_0 = a_0\) since \(\dim_{K_0}(\lambda a_0) = \dim_{K_0}a_0\) . Now the vector subspace of \(g\) generated by \(a_0\) is a non-zero ideal of \(g\) and hence is the whole of \(g\) . Hence \(g = a_0\) , which proves our assertion.

